% Part: first-order-logic
% Chapter: axiomatic-deduction
% Section: proving-things

\documentclass[../../../include/open-logic-section]{subfiles}

\begin{document}

\iftag{FOL}
      {\olfileid{fol}{axd}{pro}}
      {\olfileid{pl}{axd}{pro}}

\olsection{\usetoken{P}{derivation}નાં ઉદાહરણો}


\begin{ex}
ધારો કે આપણે $(\lnot !D \lor !E) \lif (!D \lif !E)$ સાબિત કરવા
માગીએ છીએ. સ્પષ્ટ છે કે આ આપણી કોઈ સ્વયંસિદ્ધિનો દૃષ્ટાંત નથી, તેથી
તેને !!{derive} કરવા આપણે \MP{} નિયમ વાપરવો પડશે. આપણો એકમાત્ર નિયમ
MP છે; $!A$ અને $!A \lif !B$ આપેલાં હોય ત્યારે તે~$!B$ને પ્રમાણિત
કરવા દે છે. એક યુક્તિ એ છે કે \olref[prp]{ax:lor3}માં $!A$ તરીકે
$\lnot !D$, $!B$ તરીકે $!E$ અને $!C$ તરીકે $!D \lif !E$ લઈએ, એટલે કે
નીચેનો દૃષ્ટાંત લઈએ:
\[
(\lnot !D \lif (!D \lif !E)) \lif ((!E \lif (!D \lif !E)) \lif ((\lnot
!D \lor !E) \lif (!D \lif !E))).
\]
શા માટે? MPના બે ઉપયોગથી છેલ્લો ભાગ મળે છે અને આપણને એ જ જોઈએ છે.
વળી, $\lnot !D \lif (!D \lif !E)$ એ \olref[prp]{ax:lnot2}નો અને
$!E \lif (!D \lif !E)$ એ \olref[prp]{ax:lif1}નો દૃષ્ટાંત છે, તે
સરળતાથી દેખાય છે. તેથી આપણી !!{derivation} આ છે:
\begin{derivation}
  1. &  $\lnot !D \lif (!D \lif !E)$ & \olref[prp]{ax:lnot2} \\
  2. & $(\lnot !D \lif (!D \lif !E)) \lif {}$\\
  &\qquad $((!E \lif (!D \lif !E)) \lif ((\lnot !D \lor !E) \lif (!D \lif !E)))$ & \olref[prp]{ax:lor3}\\
  3. & $(!E \lif (!D \lif !E)) \lif ((\lnot !D \lor !E) \lif (!D \lif !E))$ & 1, 2, \MP\\
  4. & $!E \lif (!D \lif !E)$ & \olref[prp]{ax:lif1}\\
  5. & $(\lnot !D \lor !E) \lif (!D \lif !E)$ & 3, 4, \MP
\end{derivation}
\end{ex}

\begin{ex}\ollabel{ex:identity}
$!D \lif !D$ની !!a{derivation} શોધવાનો પ્રયત્ન કરીએ. તે કોઈ
સ્વયંસિદ્ધિનો દૃષ્ટાંત નથી, તેથી તેને !!{derive} કરવા \MP{} વાપરવો
પડશે. \olref[prp]{ax:lif1} એ~$!A \lif !B$ સ્વરૂપની એવી સ્વયંસિદ્ધિ
છે જેના પર~\MP લાગુ કરી શકાય. તે ઉપયોગી થાય તે માટે આ કિસ્સામાં \MP{}
જે $!B$ને યોગ્ય પગલા તરીકે પ્રમાણિત કરે તે~$!D \lif !D$ જ હોવું જોઈએ,
કારણ કે આપણે આ જ !!{derive} કરવા માગીએ છીએ. તેથી $!A$ પણ $!D$ જ હોવું
જોઈએ; એટલે આપણે \olref[prp]{ax:lif1}નો નીચેનો દૃષ્ટાંત જોઈ શકીએ:
\[
!D \lif (!D \lif !D)
\]
\MP લાગુ કરવા માટે તેની અનુરૂપ બીજી આધારવિધિ, એટલે~$!A$, પણ પ્રમાણિત
કરવી પડે. આપણા કિસ્સામાં તે~$!D$ થાય, અને એકલા~$!D$ને આપણે
!!{derive} કરી શકીશું નહિ. તેથી બીજી યુક્તિ જોઈએ.

માત્ર~$\lif$ ધરાવતી બીજી સ્વયંસિદ્ધિ \olref[prp]{ax:lif2} છે, એટલે કે
\[
(!A \lif (!B \lif !C)) \lif ((!A \lif !B) \lif (!A \lif !C))
\]
\MP{}ને બે વાર લાગુ કરીને આપણે છેલ્લા અંતઃસ્થિત પ્રેરણ સુધી પહોંચી શકીએ.
ફરી તેનો અર્થ એ થાય કે આપણને \olref[prp]{ax:lif2}નો એવો દૃષ્ટાંત જોઈએ
જેમાં $!A \lif !C$ એ આપણું લક્ષ્ય !!{formula}~$!D \lif !D$ હોય. તેથી
સ્વાભાવિક રીતે $!A$ અને $!C$ બંને~$!D$ છે. હવે~$!B$ કેવી રીતે પસંદ
કરીએ જેથી $!A \lif (!B \lif !C)$ અને $!A \lif !B$, એટલે આપણા
કિસ્સામાં $!D \lif (!B \lif !D)$ અને $!D \lif !B$, બંને પણ
!!{derivable} હોય? $!B$ જે પણ લઈએ, તેમાંનું પહેલું સૂત્ર પહેલેથી જ
\olref[prp]{ax:lif1}નો દૃષ્ટાંત છે. અને જો $!B$ને $(!D \lif !D)$ લઈએ,
તો $!D \lif !B$ પણ \olref[prp]{ax:lif1}નો દૃષ્ટાંત થાય. તેથી આપણી
!!{derivation} આ છે:
\begin{derivation}
  1. & $!D \lif ((!D \lif !D) \lif !D)$ & \olref[prp]{ax:lif1}\\
  2. & $(!D \lif ((!D \lif !D) \lif !D)) \lif {}$\\
  & \qquad $((!D \lif (!D \lif !D)) \lif (!D \lif !D))$ & \olref[prp]{ax:lif2}\\
  3. & $(!D \lif (!D \lif !D)) \lif (!D \lif !D)$ & 1, 2, \MP\\
  4. & $!D \lif (!D \lif !D)$ & \olref[prp]{ax:lif1}\\
  5. & $!D \lif !D$ & 3, 4, \MP
\end{derivation}
\end{ex}

\begin{ex}\ollabel{ex:chain}
ક્યારેક આપણે બતાવવા માગીએ છીએ કે બીજા અમુક !!{formula}sના ગણ~$\Gamma$
માંથી કોઈ !!{formula}ની !!a{derivation} છે. દાખલા તરીકે, બતાવીએ કે
$\Gamma = \{!A \lif !B, !B \lif !C\}$માંથી $!A \lif !C$ને
!!{derive} કરી શકાય છે.
\begin{derivation}
  1. & $!A \lif !B$ & \Hyp\\
  2. & $!B \lif !C$ & \Hyp\\
  3. & $(!B \lif !C) \lif (!A \lif (!B \lif !C))$ & \olref[prp]{ax:lif1} \\
  4. & $!A \lif (!B \lif !C)$ & 2, 3, \MP\\
  5. & $(!A \lif (!B \lif !C)) \lif {}$\\
  & \qquad $((!A \lif !B) \lif (!A \lif !C))$ & \olref[prp]{ax:lif2}\\
  6. & $((!A \lif !B) \lif (!A \lif !C))$ & 4, 5, \MP\\
  7. & $!A \lif !C$ & 1, 6, \MP
\end{derivation}
``\Hyp'' (``પૂર્વધારણા'' માટે) લખેલી પંક્તિઓ દર્શાવે છે કે તે પંક્તિ
પરનું !!{formula} એ~$\Gamma$નું !!a{element} છે.
\end{ex}

\begin{prop}
  \ollabel{prop:chain} જો $\Gamma \Proves !A \lif !B$ અને $\Gamma
  \Proves !B \lif !C$, તો $\Gamma \Proves !A \lif !C$.
\end{prop}

\begin{proof}
  ધારો કે $\Gamma \Proves !A \lif !B$ અને $\Gamma \Proves !B \lif
  !C$. તેથી~$\Gamma$માંથી $!A \lif !B$ની !!a{derivation} છે અને
  $!B \lif !C$ની પણ~$\Gamma$માંથી !!a{derivation} છે. બંનેને એક પછી એક જોડીને એક
  !!{derivation} બનાવો. હવે અગાઉના ઉદાહરણની !!{derivation}ની પંક્તિઓ
  3--7 ઉમેરો. આ~$\Gamma$માંથી $!A \lif !C$ની !!a{derivation} છે;
  આ સૂત્ર નવી !!{derivation}ની છેલ્લી પંક્તિ છે. ધ્યાન રાખો કે
  પંક્તિ 4 અને~7નાં પ્રમાણન ત્યારે પણ માન્ય રહે છે જ્યારે પંક્તિ~1નો
  સંદર્ભ $!A \lif !B$ની !!{derivation}ની છેલ્લી પંક્તિના સંદર્ભથી અને
  પંક્તિ~2નો સંદર્ભ $!B \lif !C$ની !!{derivation}ની છેલ્લી પંક્તિના
  સંદર્ભથી બદલી દઈએ.
\end{proof}

\begin{prob}
સ્વયંસિદ્ધિઓમાંથી !!{derivation}s રજૂ કરીને બતાવો કે નીચેનું સાચું છે:
\begin{enumerate}
\item $(!A \land !B) \lif (!B \land !A)$
\item $((!A \land !B) \lif !C) \lif (!A \lif (!B \lif !C))$
\item $\lnot(!A \lor !B) \lif \lnot !A$
\end{enumerate}
\end{prob}



\end{document}
